% ============================================================
% REPORTE DE VALIDACIÓN DE HIPÓTESIS - TienditaCampus
% Autor: Emilio Jaras Sánchez
% Fecha: 26 de Marzo de 2026
% ============================================================
\documentclass[12pt, a4paper]{article}

% --- Paquetes ---
\usepackage[utf8]{inputenc}
\usepackage[T1]{fontenc}
\usepackage[spanish]{babel}
\usepackage{geometry}
\usepackage{graphicx}
\usepackage{booktabs}
\usepackage{array}
\usepackage{longtable}
\usepackage{xcolor}
\usepackage{colortbl}
\usepackage{hyperref}
\usepackage{fancyhdr}
\usepackage{titlesec}
\usepackage{amsmath}
\usepackage{float}
\usepackage{multirow}
\usepackage{caption}
\usepackage{subcaption}
\usepackage{pgfplots}
\pgfplotsset{compat=1.18}
\usepackage{tikz}

\geometry{margin=2.5cm}

% --- Colores Corporativos ---
\definecolor{primaryblue}{HTML}{1E3A5F}
\definecolor{accentgreen}{HTML}{2E7D32}
\definecolor{warnorange}{HTML}{E65100}
\definecolor{lightgray}{HTML}{F5F5F5}
\definecolor{tableheader}{HTML}{1E3A5F}
\definecolor{rowalt}{HTML}{E8EAF6}
\definecolor{successgreen}{HTML}{4CAF50}
\definecolor{dangered}{HTML}{F44336}

% --- Estilo de Headers ---
\pagestyle{fancy}
\fancyhf{}
\fancyhead[L]{\small\textcolor{primaryblue}{\textbf{TienditaCampus}}}
\fancyhead[R]{\small\textcolor{primaryblue}{Reporte de Validación de Hipótesis}}
\fancyfoot[C]{\thepage}
\renewcommand{\headrulewidth}{0.4pt}

% --- Estilo de Títulos ---
\titleformat{\section}{\Large\bfseries\color{primaryblue}}{\thesection}{1em}{}[\titlerule]
\titleformat{\subsection}{\large\bfseries\color{primaryblue!80}}{\thesubsection}{1em}{}

% --- Hyperref ---
\hypersetup{
    colorlinks=true,
    linkcolor=primaryblue,
    urlcolor=primaryblue,
    citecolor=primaryblue
}

\begin{document}

% ============================================================
% PORTADA
% ============================================================
\begin{titlepage}
    \centering
    \vspace*{2cm}
    
    % Logo placeholder
    \begin{tikzpicture}
        \draw[fill=primaryblue, rounded corners=8pt] (0,0) rectangle (3,3);
        \node[white, font=\Huge\bfseries] at (1.5,1.5) {TC};
    \end{tikzpicture}
    
    \vspace{1.5cm}
    
    {\Huge\bfseries\textcolor{primaryblue}{Reporte de Validación\\[0.3cm] de Hipótesis}}
    
    \vspace{0.8cm}
    
    {\Large\textcolor{primaryblue!70}{Análisis de Datos para la Evaluación de\\Sostenibilidad del Marketplace Universitario}}
    
    \vspace{1.5cm}
    
    \begin{tikzpicture}
        \draw[primaryblue, thick] (0,0) -- (12,0);
    \end{tikzpicture}
    
    \vspace{1.5cm}
    
    {\large\textbf{Proyecto:} TienditaCampus}\\[0.3cm]
    {\large\textbf{Autor:} Emilio Jaras Sánchez}\\[0.3cm]
    {\large\textbf{Correo:} jarassabchezl@gmail.com}\\[0.3cm]
    {\large\textbf{Fecha:} 26 de Marzo de 2026}\\[0.3cm]
    {\large\textbf{Periodo de Datos:} 9 -- 20 de Marzo de 2026}
    
    \vspace{1.5cm}
    
    \begin{tikzpicture}
        \draw[primaryblue, thick] (0,0) -- (12,0);
    \end{tikzpicture}
    
    \vspace{1cm}
    
    {\normalsize\textit{Documento generado a partir de datos reales\\extraídos de la base de datos PostgreSQL de producción.}}
    
\end{titlepage}

% ============================================================
% ÍNDICE
% ============================================================
\tableofcontents
\newpage

% ============================================================
% 1. INTRODUCCIÓN
% ============================================================
\section{Introducción}

\subsection{Contexto del Proyecto}
\textbf{TienditaCampus} es una plataforma de marketplace universitario diseñada para conectar vendedores estudiantiles con compradores dentro del campus. El proyecto busca demostrar que un ecosistema de microcomercio digital es viable, rentable y sostenible dentro de una comunidad universitaria.

\subsection{Hipótesis a Validar}

\begin{center}
\fcolorbox{primaryblue}{lightgray}{%
    \parbox{0.9\textwidth}{%
        \textbf{\textcolor{primaryblue}{Hipótesis (H1):}} ``La implementación de una plataforma digital de marketplace universitario permite a los vendedores estudiantiles alcanzar márgenes de ganancia superiores al 20\%, con una tasa de pérdidas (merma) inferior al 15\%, demostrando la sostenibilidad económica del modelo de negocio en un periodo operativo de 2 semanas.''
    }%
}
\end{center}

\vspace{0.5cm}

\begin{center}
\fcolorbox{warnorange}{lightgray}{%
    \parbox{0.9\textwidth}{%
        \textbf{\textcolor{warnorange}{Hipótesis Nula (H0):}} ``El marketplace universitario no logra márgenes de ganancia superiores al 20\% ni una tasa de pérdidas inferior al 15\%, por lo que el modelo no es económicamente sostenible.''
    }%
}
\end{center}

\subsection{Criterios de Validación}

\begin{table}[H]
\centering
\caption{Criterios de aceptación para la validación de hipótesis}
\begin{tabular}{>{\bfseries}l c c}
\toprule
\rowcolor{tableheader}
\textcolor{white}{\textbf{Métrica}} & \textcolor{white}{\textbf{Umbral Mínimo}} & \textcolor{white}{\textbf{Meta Óptima}} \\
\midrule
Margen de Ganancia Promedio & $\geq 20\%$ & $\geq 40\%$ \\
\rowcolor{rowalt}
Tasa de Pérdidas (Merma) & $\leq 15\%$ & $\leq 5\%$ \\
Días Rentables & $\geq 60\%$ & $\geq 80\%$ \\
\rowcolor{rowalt}
Punto de Equilibrio & Alcanzado & Superado \\
Compradores Activos & $\geq 10$ & $\geq 25$ \\
\bottomrule
\end{tabular}
\end{table}

% ============================================================
% 2. METODOLOGÍA
% ============================================================
\section{Metodología}

\subsection{Fuente de Datos}
Los datos provienen de la base de datos PostgreSQL de producción de TienditaCampus, alojada en un servidor AWS EC2 (IP: 3.219.19.219). Se utilizaron las siguientes tablas:

\begin{itemize}
    \item \texttt{users} -- Registro de vendedores y compradores
    \item \texttt{products} -- Catálogo de productos disponibles
    \item \texttt{categories} -- Clasificación de productos
    \item \texttt{orders} / \texttt{order\_items} -- Transacciones de compra-venta
    \item \texttt{daily\_sales} / \texttt{sale\_details} -- Métricas diarias de venta
\end{itemize}

\subsection{Periodo de Análisis}
\begin{itemize}
    \item \textbf{Semana 1:} 9 al 13 de Marzo de 2026 (5 días hábiles)
    \item \textbf{Semana 2:} 17 al 20 de Marzo de 2026 (4 días hábiles)
    \item \textbf{Total:} 9 días operativos
\end{itemize}

\subsection{Herramientas de Análisis}
\begin{itemize}
    \item \textbf{Base de Datos:} PostgreSQL 15 (Docker)
    \item \textbf{Backend:} NestJS con TypeORM
    \item \textbf{Dashboard:} Next.js 14 con gráficas Recharts
    \item \textbf{BigQuery:} Snapshots para benchmarking externo
\end{itemize}

% ============================================================
% 3. DATOS DEL ECOSISTEMA
% ============================================================
\section{Datos del Ecosistema}

\subsection{Participantes del Marketplace}

\begin{table}[H]
\centering
\caption{Distribución de usuarios en el ecosistema}
\begin{tabular}{l c c}
\toprule
\rowcolor{tableheader}
\textcolor{white}{\textbf{Rol}} & \textcolor{white}{\textbf{Cantidad}} & \textcolor{white}{\textbf{\% del Total}} \\
\midrule
Vendedores (Sellers) & 5 & 12.5\% \\
\rowcolor{rowalt}
Compradores (Buyers) & 35 & 87.5\% \\
\midrule
\textbf{Total Usuarios} & \textbf{40} & \textbf{100\%} \\
\bottomrule
\end{tabular}
\end{table}

\subsection{Vendedores Registrados}

\begin{table}[H]
\centering
\caption{Detalle de los 5 vendedores activos del campus}
\begin{tabular}{c l l l}
\toprule
\rowcolor{tableheader}
\textcolor{white}{\textbf{\#}} & \textcolor{white}{\textbf{Nombre}} & \textcolor{white}{\textbf{Email}} & \textcolor{white}{\textbf{Productos}} \\
\midrule
1 & Ana García & anagarpep14@gmail.com & 4 \\
\rowcolor{rowalt}
2 & Diego Ramírez & diegoramiraasa2@gmail.com & 4 \\
3 & Elena Gómez & elenaaaszgomezr@gmail.com & 4 \\
\rowcolor{rowalt}
4 & Carlos Pérez & carlospereeezagui23@gmail.com & 4 \\
5 & Sofía Lara & sofialaaaaaar2535@gmail.com & 4 \\
\bottomrule
\end{tabular}
\end{table}

\subsection{Catálogo de Productos}

\begin{table}[H]
\centering
\caption{Productos ofertados en el marketplace con análisis de márgenes}
\small
\begin{tabular}{l l r r r c}
\toprule
\rowcolor{tableheader}
\textcolor{white}{\textbf{Producto}} & \textcolor{white}{\textbf{Categoría}} & \textcolor{white}{\textbf{Costo}} & \textcolor{white}{\textbf{Precio}} & \textcolor{white}{\textbf{Margen}} & \textcolor{white}{\textbf{Perec.}} \\
\midrule
Hamburguesa & Comida Preparada & \$15.38 & \$32.58 & 52.8\% & Sí \\
\rowcolor{rowalt}
Chicharrines & Comida Preparada & \$18.67 & \$39.55 & 52.8\% & Sí \\
Burritos & Snacks & \$10.16 & \$21.52 & 52.8\% & No \\
\rowcolor{rowalt}
Chicles & Snacks & \$5.61 & \$11.88 & 52.8\% & No \\
Burrito (Diego) & Comida Preparada & \$16.41 & \$27.64 & 40.6\% & Sí \\
\rowcolor{rowalt}
Ensalada & Comida Preparada & \$22.58 & \$38.04 & 40.6\% & Sí \\
Papitas & Snacks & \$8.93 & \$15.04 & 40.6\% & No \\
\rowcolor{rowalt}
Papas adobadas & Snacks & \$7.94 & \$13.38 & 40.7\% & No \\
Ensalada César & Comida Preparada & \$27.14 & \$56.29 & 51.8\% & Sí \\
\rowcolor{rowalt}
Sandwich & Comida Preparada & \$26.38 & \$54.71 & 51.8\% & Sí \\
Galletas Chokis & Snacks & \$9.89 & \$20.51 & 51.8\% & No \\
\rowcolor{rowalt}
Coca Cola & Bebidas & \$6.95 & \$14.41 & 51.8\% & No \\
Carlotas & Comida Preparada & \$20.05 & \$24.35 & 17.7\% & Sí \\
\rowcolor{rowalt}
Cheesecake & Comida Preparada & \$20.70 & \$25.14 & 17.7\% & Sí \\
Platanitos & Snacks & \$10.93 & \$13.27 & 17.6\% & No \\
\rowcolor{rowalt}
Paletas & Bebidas & \$6.85 & \$8.32 & 17.7\% & No \\
Burrito Pollo & Comida Preparada & \$23.04 & \$27.38 & 15.8\% & Sí \\
\rowcolor{rowalt}
Sandwich de Atún & Comida Preparada & \$16.18 & \$19.22 & 15.8\% & Sí \\
Barritas & Snacks & \$10.81 & \$12.84 & 15.8\% & No \\
\rowcolor{rowalt}
Galletas Oreo & Snacks & \$11.20 & \$13.31 & 15.8\% & No \\
\bottomrule
\end{tabular}
\end{table}

\subsection{Distribución por Categoría}

\begin{table}[H]
\centering
\caption{Productos por categoría}
\begin{tabular}{l c c}
\toprule
\rowcolor{tableheader}
\textcolor{white}{\textbf{Categoría}} & \textcolor{white}{\textbf{Productos}} & \textcolor{white}{\textbf{\% del Catálogo}} \\
\midrule
Comida Preparada & 10 & 50\% \\
\rowcolor{rowalt}
Snacks & 8 & 40\% \\
Bebidas & 2 & 10\% \\
\midrule
\textbf{Total} & \textbf{20} & \textbf{100\%} \\
\bottomrule
\end{tabular}
\end{table}

% ============================================================
% 4. ANÁLISIS DE VENTAS
% ============================================================
\section{Análisis de Ventas}

\subsection{Resumen General de Transacciones}

A partir de los registros de la tabla \texttt{orders}, se identificaron las siguientes métricas globales:

\begin{table}[H]
\centering
\caption{Métricas globales de transacciones}
\begin{tabular}{l r}
\toprule
\rowcolor{tableheader}
\textcolor{white}{\textbf{Métrica}} & \textcolor{white}{\textbf{Valor}} \\
\midrule
Total de Órdenes Completadas & 48 \\
\rowcolor{rowalt}
Total de Artículos Vendidos (order\_items) & 98 \\
Compradores Activos (únicos) & 27 \\
\rowcolor{rowalt}
Monto Total de Ventas & \$4,408.24 MXN \\
Ticket Promedio por Orden & \$91.84 MXN \\
\rowcolor{rowalt}
Artículos Promedio por Orden & 2.04 \\
\bottomrule
\end{tabular}
\end{table}

\subsection{Ventas por Vendedor}

\begin{table}[H]
\centering
\caption{Desempeño de ventas por vendedor durante el periodo de análisis}
\begin{tabular}{l c r r r}
\toprule
\rowcolor{tableheader}
\textcolor{white}{\textbf{Vendedor}} & \textcolor{white}{\textbf{Órdenes}} & \textcolor{white}{\textbf{Ingresos}} & \textcolor{white}{\textbf{Ticket Prom.}} & \textcolor{white}{\textbf{Margen Est.}} \\
\midrule
Ana García & 10 & \$961.85 & \$96.19 & $\sim$52.8\% \\
\rowcolor{rowalt}
Diego Ramírez & 11 & \$1,308.24 & \$118.93 & $\sim$40.6\% \\
Elena Gómez & 8 & \$1,002.64 & \$125.33 & $\sim$51.8\% \\
\rowcolor{rowalt}
Carlos Pérez & 10 & \$691.57 & \$69.16 & $\sim$17.7\% \\
Sofía Lara & 9 & \$443.94 & \$49.33 & $\sim$15.8\% \\
\midrule
\textbf{Total} & \textbf{48} & \textbf{\$4,408.24} & \textbf{\$91.84} & -- \\
\bottomrule
\end{tabular}
\end{table}

\subsection{Ventas por Semana}

\begin{table}[H]
\centering
\caption{Comparativa semanal de actividad comercial}
\begin{tabular}{l c r r}
\toprule
\rowcolor{tableheader}
\textcolor{white}{\textbf{Semana}} & \textcolor{white}{\textbf{Órdenes}} & \textcolor{white}{\textbf{Ingresos}} & \textcolor{white}{\textbf{Promedio Diario}} \\
\midrule
Semana 1 (Mar 9--13) & 26 & \$2,467.71 & \$493.54 \\
\rowcolor{rowalt}
Semana 2 (Mar 17--20) & 22 & \$1,940.53 & \$485.13 \\
\midrule
\textbf{Total} & \textbf{48} & \textbf{\$4,408.24} & \textbf{\$489.80} \\
\bottomrule
\end{tabular}
\end{table}

\subsection{Top 10 Productos Más Vendidos (por unidades)}

\begin{table}[H]
\centering
\caption{Productos con mayor volumen de ventas}
\begin{tabular}{c l l c r}
\toprule
\rowcolor{tableheader}
\textcolor{white}{\textbf{\#}} & \textcolor{white}{\textbf{Producto}} & \textcolor{white}{\textbf{Categoría}} & \textcolor{white}{\textbf{Uds.}} & \textcolor{white}{\textbf{Ingreso}} \\
\midrule
1 & Ensalada & Comida Prep. & 16 & \$608.64 \\
\rowcolor{rowalt}
2 & Burrito (Diego) & Comida Prep. & 14 & \$386.96 \\
3 & Cheesecake & Comida Prep. & 13 & \$326.82 \\
\rowcolor{rowalt}
4 & Chicharrines & Comida Prep. & 11 & \$435.05 \\
5 & Platanitos & Snacks & 11 & \$145.97 \\
\rowcolor{rowalt}
6 & Galletas Oreo & Snacks & 10 & \$133.10 \\
7 & Burritos & Snacks & 9 & \$193.68 \\
\rowcolor{rowalt}
8 & Burrito Pollo & Comida Prep. & 10 & \$273.80 \\
9 & Sandwich & Comida Prep. & 8 & \$437.68 \\
\rowcolor{rowalt}
10 & Papitas & Snacks & 9 & \$135.36 \\
\bottomrule
\end{tabular}
\end{table}

% ============================================================
% 5. ANÁLISIS DE RENTABILIDAD
% ============================================================
\section{Análisis de Rentabilidad}

\subsection{Margen de Ganancia por Vendedor}

\begin{table}[H]
\centering
\caption{Análisis de rentabilidad por vendedor}
\begin{tabular}{l r r r c}
\toprule
\rowcolor{tableheader}
\textcolor{white}{\textbf{Vendedor}} & \textcolor{white}{\textbf{Ingreso}} & \textcolor{white}{\textbf{Costo Est.}} & \textcolor{white}{\textbf{Ganancia}} & \textcolor{white}{\textbf{Margen}} \\
\midrule
Ana García & \$961.85 & \$453.67 & \$508.18 & \textcolor{successgreen}{\textbf{52.8\%}} \\
\rowcolor{rowalt}
Diego Ramírez & \$1,308.24 & \$777.10 & \$531.14 & \textcolor{successgreen}{\textbf{40.6\%}} \\
Elena Gómez & \$1,002.64 & \$483.27 & \$519.37 & \textcolor{successgreen}{\textbf{51.8\%}} \\
\rowcolor{rowalt}
Carlos Pérez & \$691.57 & \$569.37 & \$122.20 & \textcolor{warnorange}{\textbf{17.7\%}} \\
Sofía Lara & \$443.94 & \$373.80 & \$70.14 & \textcolor{dangered}{\textbf{15.8\%}} \\
\midrule
\textbf{Promedio Pond.} & \textbf{\$4,408.24} & \textbf{\$2,657.21} & \textbf{\$1,751.03} & \textbf{39.7\%} \\
\bottomrule
\end{tabular}
\end{table}

\subsection{Interpretación de Márgenes}

\begin{itemize}
    \item \textcolor{successgreen}{\textbf{3 de 5 vendedores}} (60\%) superan el umbral mínimo de 20\% de margen con holgura (40--53\%).
    \item \textcolor{warnorange}{\textbf{1 vendedor}} (Carlos Pérez, 17.7\%) está \textit{cerca} del umbral pero sin alcanzarlo -- requiere ajuste de precios.
    \item \textcolor{dangered}{\textbf{1 vendedor}} (Sofía Lara, 15.8\%) opera con márgenes insuficientes.
    \item El \textbf{margen promedio ponderado global} es de \textbf{39.7\%}, lo que \textbf{supera ampliamente} el umbral de 20\%.
\end{itemize}

\subsection{Gráfica de Márgenes por Vendedor}

\begin{figure}[H]
\centering
\begin{tikzpicture}
\begin{axis}[
    ybar,
    bar width=20pt,
    width=0.9\textwidth,
    height=7cm,
    ylabel={Margen de Ganancia (\%)},
    symbolic x coords={Ana, Diego, Elena, Carlos, Sofía},
    xtick=data,
    ymin=0, ymax=60,
    nodes near coords,
    nodes near coords align={vertical},
    every node near coord/.append style={font=\small\bfseries},
    axis lines*=left,
    grid=major,
    grid style={dashed, gray!30},
    extra y ticks={20},
    extra y tick style={grid style={red, thick, dashed}},
    extra y tick labels={Umbral 20\%},
]
\addplot[fill=primaryblue!80] coordinates {
    (Ana, 52.8) (Diego, 40.6) (Elena, 51.8) (Carlos, 17.7) (Sofía, 15.8)
};
\end{axis}
\end{tikzpicture}
\caption{Margen de ganancia por vendedor vs. umbral mínimo de 20\%}
\end{figure}

% ============================================================
% 6. ANÁLISIS DE PÉRDIDAS (MERMA)
% ============================================================
\section{Análisis de Pérdidas y Merma}

\subsection{Datos de Merma por Categoría}

Basándonos en los registros de \texttt{daily\_sales} y \texttt{sale\_details}, se estima la merma considerando la proporción de productos perecederos:

\begin{table}[H]
\centering
\caption{Estimación de pérdidas por categoría de producto}
\begin{tabular}{l c c r c}
\toprule
\rowcolor{tableheader}
\textcolor{white}{\textbf{Categoría}} & \textcolor{white}{\textbf{Uds. Vendidas}} & \textcolor{white}{\textbf{Uds. Perdidas}} & \textcolor{white}{\textbf{Costo Merma}} & \textcolor{white}{\textbf{Tasa}} \\
\midrule
Comida Preparada & 82 & 8 & \$164.40 & \textcolor{successgreen}{8.9\%} \\
\rowcolor{rowalt}
Snacks & 58 & 2 & \$19.60 & \textcolor{successgreen}{3.3\%} \\
Bebidas & 6 & 0 & \$0.00 & \textcolor{successgreen}{0.0\%} \\
\midrule
\textbf{Total} & \textbf{146} & \textbf{10} & \textbf{\$184.00} & \textbf{6.4\%} \\
\bottomrule
\end{tabular}
\end{table}

\subsection{Evaluación de la Tasa de Pérdidas}

\begin{center}
\fcolorbox{accentgreen}{lightgray}{%
    \parbox{0.85\textwidth}{%
        \textbf{\textcolor{accentgreen}{Resultado:}} La tasa de pérdidas global estimada es de \textbf{6.4\%}, lo cual se encuentra \textbf{por debajo} del umbral máximo de 15\% y se clasifica como \textbf{``Bueno''} según los criterios de evaluación.
    }%
}
\end{center}

\begin{itemize}
    \item $\leq 5\%$: \textcolor{successgreen}{\textbf{EXCELENTE}} -- Meta de industria
    \item $5\% - 10\%$: \textcolor{accentgreen}{\textbf{BUENO}} -- Aceptable para microcomercio $\leftarrow$ \textbf{Posición actual (6.4\%)}
    \item $10\% - 15\%$: \textcolor{warnorange}{\textbf{ACEPTABLE}} -- Requiere optimización
    \item $> 15\%$: \textcolor{dangered}{\textbf{CRÍTICO}} -- Modelo insostenible
\end{itemize}

% ============================================================
% 7. ANÁLISIS DE ADOPCIÓN
% ============================================================
\section{Análisis de Adopción de la Plataforma}

\subsection{Métricas de Participación}

\begin{table}[H]
\centering
\caption{Indicadores de adopción del marketplace}
\begin{tabular}{l r c}
\toprule
\rowcolor{tableheader}
\textcolor{white}{\textbf{Indicador}} & \textcolor{white}{\textbf{Valor}} & \textcolor{white}{\textbf{Evaluación}} \\
\midrule
Compradores registrados & 35 & \textcolor{successgreen}{\textbf{Supera meta}} \\
\rowcolor{rowalt}
Compradores activos (con compras) & 27 & \textcolor{successgreen}{\textbf{Supera meta}} \\
Tasa de activación & 77.1\% & \textcolor{successgreen}{\textbf{Excelente}} \\
\rowcolor{rowalt}
Vendedores activos & 5 & \textcolor{successgreen}{\textbf{100\% activos}} \\
Promedio compras por buyer & 1.78 & \textcolor{accentgreen}{\textbf{Bueno}} \\
\rowcolor{rowalt}
Días operativos & 9 & -- \\
\bottomrule
\end{tabular}
\end{table}

\subsection{Gráfica de Distribución de Órdenes}

\begin{figure}[H]
\centering
\begin{tikzpicture}
\begin{axis}[
    ybar,
    bar width=15pt,
    width=0.9\textwidth,
    height=7cm,
    ylabel={Número de Órdenes},
    xlabel={Día de operación},
    symbolic x coords={Mar 9, Mar 10, Mar 11, Mar 12, Mar 13, Mar 17, Mar 18, Mar 19, Mar 20},
    xtick=data,
    x tick label style={rotate=45, anchor=east, font=\small},
    ymin=0, ymax=10,
    nodes near coords,
    every node near coord/.append style={font=\tiny\bfseries},
    axis lines*=left,
    grid=major,
    grid style={dashed, gray!30},
]
\addplot[fill=primaryblue!70] coordinates {
    (Mar 9, 6) (Mar 10, 6) (Mar 11, 6) (Mar 12, 4) (Mar 13, 4) (Mar 17, 5) (Mar 18, 5) (Mar 19, 5) (Mar 20, 7)
};
\end{axis}
\end{tikzpicture}
\caption{Distribución diaria de órdenes completadas}
\end{figure}

% ============================================================
% 8. ANÁLISIS DE SOSTENIBILIDAD
% ============================================================
\section{Análisis de Sostenibilidad del Modelo}

\subsection{Indicadores Clave de Desempeño (KPIs)}

\begin{table}[H]
\centering
\caption{Dashboard de KPIs -- Evaluación de Sostenibilidad}
\begin{tabular}{l r l c}
\toprule
\rowcolor{tableheader}
\textcolor{white}{\textbf{Métrica}} & \textcolor{white}{\textbf{Valor}} & \textcolor{white}{\textbf{Unidad}} & \textcolor{white}{\textbf{Evaluación}} \\
\midrule
Margen Promedio (ponderado) & 39.7 & \% & \textcolor{successgreen}{\textbf{EXCELENTE}} \\
\rowcolor{rowalt}
Tasa de Pérdidas & 6.4 & \% & \textcolor{accentgreen}{\textbf{BUENO}} \\
Días Rentables & 9 de 9 & días & \textcolor{successgreen}{\textbf{EXCELENTE}} \\
\rowcolor{rowalt}
Ingreso Diario Promedio & 489.80 & \$/día & \textcolor{successgreen}{\textbf{BUENO}} \\
Ganancia Total Estimada & 1,751.03 & \$ & \textcolor{successgreen}{\textbf{POSITIVA}} \\
\rowcolor{rowalt}
Compradores Únicos & 27 & usuarios & \textcolor{successgreen}{\textbf{SUPERA META}} \\
\bottomrule
\end{tabular}
\end{table}

\subsection{Proyección a 30 Días}

\begin{table}[H]
\centering
\caption{Proyección financiera a 30 días operativos}
\begin{tabular}{l r r}
\toprule
\rowcolor{tableheader}
\textcolor{white}{\textbf{Concepto}} & \textcolor{white}{\textbf{Periodo Actual (9 días)}} & \textcolor{white}{\textbf{Proyección (30 días)}} \\
\midrule
Ingresos Totales & \$4,408.24 & \$14,694.13 \\
\rowcolor{rowalt}
Costos Estimados & \$2,657.21 & \$8,857.37 \\
Ganancia Bruta & \$1,751.03 & \$5,836.77 \\
\rowcolor{rowalt}
Pérdida por Merma & \$184.00 & \$613.33 \\
\textbf{Ganancia Neta} & \textbf{\$1,567.03} & \textbf{\$5,223.43} \\
\bottomrule
\end{tabular}
\end{table}

% ============================================================
% 9. VALIDACIÓN DE HIPÓTESIS
% ============================================================
\section{Validación de Hipótesis}

\subsection{Matriz de Evaluación Final}

\begin{table}[H]
\centering
\caption{\textbf{Matriz de validación: resultados vs. criterios de aceptación}}
\begin{tabular}{l c c c}
\toprule
\rowcolor{tableheader}
\textcolor{white}{\textbf{Criterio}} & \textcolor{white}{\textbf{Umbral}} & \textcolor{white}{\textbf{Resultado}} & \textcolor{white}{\textbf{Estado}} \\
\midrule
Margen $\geq$ 20\% & 20\% & 39.7\% & \textcolor{successgreen}{\textbf{$\checkmark$ APROBADO}} \\
\rowcolor{rowalt}
Tasa pérdida $\leq$ 15\% & 15\% & 6.4\% & \textcolor{successgreen}{\textbf{$\checkmark$ APROBADO}} \\
Días rentables $\geq$ 60\% & 60\% & 100\% & \textcolor{successgreen}{\textbf{$\checkmark$ APROBADO}} \\
\rowcolor{rowalt}
Compradores $\geq$ 10 & 10 & 27 & \textcolor{successgreen}{\textbf{$\checkmark$ APROBADO}} \\
Punto de equilibrio & Alcanzado & Superado & \textcolor{successgreen}{\textbf{$\checkmark$ APROBADO}} \\
\midrule
\multicolumn{3}{r}{\textbf{Criterios aprobados:}} & \textbf{5 / 5} \\
\bottomrule
\end{tabular}
\end{table}

\subsection{Resultado}

\begin{center}
\begin{tikzpicture}
\draw[fill=accentgreen!15, draw=accentgreen, thick, rounded corners=10pt] (0,0) rectangle (14,3.5);
\node[font=\LARGE\bfseries, text=accentgreen] at (7,2.5) {HIPÓTESIS H1: VALIDADA};
\node[font=\normalsize, text width=12cm, align=center] at (7,1.2) {
Los datos demuestran que el marketplace universitario TienditaCampus es económicamente sostenible, con un margen de ganancia promedio del 39.7\% (superando el umbral de 20\%), una tasa de pérdidas del 6.4\% (por debajo del 15\%) y 27 compradores activos (superando la meta de 10).
};
\end{tikzpicture}
\end{center}

\vspace{0.5cm}

\textbf{Se rechaza la hipótesis nula (H0)} y se \textbf{acepta la hipótesis alternativa (H1)}: el modelo de marketplace universitario es económicamente sostenible.

% ============================================================
% 10. RECOMENDACIONES
% ============================================================
\section{Recomendaciones}

\begin{enumerate}
    \item \textbf{Optimizar márgenes de vendedores con baja rentabilidad:} Carlos Pérez (17.7\%) y Sofía Lara (15.8\%) deben revisar su estructura de costos o incrementar sus precios de venta.
    
    \item \textbf{Reducir merma en comida preparada:} Implementar un sistema de producción bajo demanda basado en los patrones de pedidos históricos, reduciendo las 8 unidades perdidas en esta categoría.
    
    \item \textbf{Expandir la categoría de bebidas:} Con solo 2 productos y 0\% de merma, las bebidas representan una oportunidad de crecimiento con riesgo mínimo.
    
    \item \textbf{Aumentar la recurrencia de compra:} El promedio de 1.78 compras/buyer sugiere oportunidad de implementar programas de lealtad o descuentos por volumen.
    
    \item \textbf{Escalar a más vendedores:} Con 5 vendedores atendiendo a 27 compradores activos, existe capacidad para incorporar más microemprendedores estudiantiles.
\end{enumerate}

% ============================================================
% 11. ANEXOS
% ============================================================
\section{Anexos}

\subsection{Anexo A: Consultas SQL Utilizadas}

Las siguientes consultas se utilizaron para la extracción de datos (disponibles en \texttt{consultas\_analisis.sql}):

\begin{enumerate}
    \item Análisis General de Ventas por Semana
    \item Rendimiento por Categoría de Producto
    \item Análisis de Rentabilidad por Producto
    \item Análisis de Pérdidas por Motivo
    \item Tendencia de Ventas Diarias
    \item Análisis de Punto de Equilibrio
    \item Resumen Semanal Consolidado
    \item Análisis de Eficiencia Operativa
    \item Top 5 Productos Más Rentables
    \item Análisis de Sostenibilidad del Negocio
\end{enumerate}

\subsection{Anexo B: Infraestructura Técnica}

\begin{table}[H]
\centering
\caption{Stack tecnológico del proyecto}
\begin{tabular}{l l}
\toprule
\rowcolor{tableheader}
\textcolor{white}{\textbf{Componente}} & \textcolor{white}{\textbf{Tecnología}} \\
\midrule
Frontend & Next.js 14, React 18, TypeScript \\
\rowcolor{rowalt}
Backend & NestJS, TypeORM, PostgreSQL 15 \\
Hosting Frontend & Vercel (rama \texttt{jesel}) \\
\rowcolor{rowalt}
Hosting Backend & AWS EC2 (Docker) \\
Base de Datos & PostgreSQL 15 (AWS EC2) \\
\rowcolor{rowalt}
Benchmarking & Google BigQuery \\
Autenticación & JWT + Google OAuth 2.0 \\
\bottomrule
\end{tabular}
\end{table}

\vspace{1cm}

\begin{center}
\begin{tikzpicture}
\draw[primaryblue, thick] (0,0) -- (12,0);
\end{tikzpicture}

\vspace{0.5cm}
\textit{Fin del Reporte -- TienditaCampus © 2026}
\end{center}

\end{document}
